\section{The smaller zero-normal companion}
\label{sec:zero}

There is a second presentation, $\mathcal E^0_3$, with the same three
inference families and the same fixed axiom clause. Its terms are canonical
only for the deletion
\[
                         I(0,a;u.h)\longrightarrow a.
\]
Let $N_0$ perform this deletion hereditarily. Every rule substitution is
followed by $N_0$. A successor-headed iterator is not unfolded by this
syntactic operation. The evaluation map is instead
\[
 \sigma_E(F)=
 \begin{cases}
  r=N_0\bigl(h[I(c,t;u.h)/u]\bigr),
    &E\equiv(s=t),\quad F\equiv(I(Sc,s;u.h)=r),\\[3pt]
  F,&\text{otherwise}.
 \end{cases}
\]
The initializer in the equation $E$ must match the initializer of the
iterator in $F$. The previous iterate is constructed from the replacement
initializer $t$; it is not an additional premise. This is precisely the
existing SeedEval rule, executed only at witness discharge.

Use and Intro are those of Section~\ref{sec:raw}, with normalized
substitution and the literal witness test on canonical syntax. Discharge
exports $\sigma_E(F)$ instead of $\rho_E(F)$, subject to the same outer
scope condition. There is one arithmetic pattern and the identity case.

The proof of Theorem~\ref{thm:raw-equivalence} applies verbatim with
SeedEval in place of the two raw Eval cases. Its separate-rule control has
Use, Intro, Open and SeedEval. The compound major still recovers Open, and
$\Theta$ still recovers SeedEval. In the reverse simulation no arbitrary
substitution into an existing derivation is assumed.

\paragraph{The arithmetic interface is unchanged.}
At counter $0$, SeedEval reads
\[
             a=b,\ I(S0,a;u.h)=r\quad\Longrightarrow\quad r=N_0(h[b/u]).
\]
Taking $h=u$ and using two such evaluations derives the joining rule for
equations. Taking arbitrary $h$ derives transport by the same two-step
pivot construction. Two reflexivities give successor computation.
Zero computation is canonical syntax. This interface is formalized in the
pinned \texttt{NativeZeroSeed.lean} as \texttt{NVPrf}; the equivalence
\texttt{nvprf\_iff\_pra} is already in that source~\cite{zeroformal}.
The clause representation and its fixed reflexivity clause are formalized
in \texttt{NativeZeroClausal.lean}. The present fusion is a new written
equivalence on top of those interfaces, not a claim that its new checker
has been verified by Lean.

\paragraph{Only structural normalization is used.}
Recursively normalize an iterator's children. If its normalized counter is
$0$, return its normalized initializer; otherwise rebuild the iterator.
Every contraction deletes nodes, and none duplicates a body. Structural
induction proves termination, idempotence and substitution coherence
\[
              N_0(N_0(t)[\tau])=N_0(t[\tau]).
\]
Thus this companion does not reinstate the full computation normalizer
removed from the main presentation.

\section{Full first-order iPRA and both interpretation composites}
\label{sec:interpretations}

The reference theory is full first-order intuitionistic primitive recursive
arithmetic: intuitionistic equality logic; a symbol for each fixed finite
primitive-recursive description and its defining equations; separation of
zero from successors; successor injectivity; and induction for every
quantifier-free formula, with parameters and retained assumptions.
In particular, the target is not merely equational PRA or a negative
fragment of iPRA.

\subsection{The intermediate ordinary scalar theory}

Let $T_I$ be ordinary intuitionistic first-order equality logic over the
staged scalar-iterator signature. Its nonlogical axioms and rules are exactly
the two iterator equations, numerical coverage, separation, and contextual
equation induction:
\begin{gather*}
 I(0,a;u.h)=a,\qquad I(Sn,a;u.h)=h[I(n,a;u.h)/u],\\
 \forall x\,(x=0\vee\exists y\,x=Sy),\qquad (0=S0)\to\bot,\\
 \frac{\Gamma\vdash E(0)\qquad\Gamma,E(z)\vdash E(Sz)}
      {\Gamma\vdash E(t)}\quad(E\text{ an equation},\ z\text{ fresh}).
\end{gather*}
The displayed separation sentence is essential to the ordinary realization:
positive equations by themselves do not prohibit a one-element structure.
It is not an extra hidden axiom of the native calculus, where $0=S0$ is
the designated contradiction and explosion is derived.

The native-to-ordinary translation is
\[
 (s=t)^\circ=(s=t),\qquad
 (D_{x,y}B)^\circ=\exists x\bigl(B^\circ(x,0)\to\forall y\,B^\circ(x,y)\bigr).
\]
Unconditional clauses translate to universals and conditional clauses to
implications into universals. All numerical domains are the whole domain,
defined by $x=x$, and equality is actual equality.

\subsection{Explicit logical macros}

Put $C(v,a,b)=I(v,a;u.b)$, with $u$ unused, and
$p(z)=I(z,0;u.C(u,1,0))$. The derived computation equations give the
selector endpoints and $p(1)=1$, $p(2)=0$.
For a formula, define its gate recursively by
\[
 G_v(s=t)=(C(v,0,s)=C(v,0,t)),\qquad
 G_v(D_{x,y}B)=D_{x,y}G_vB,
\]
with fresh bound variables. It is provable at $v=0$ and equivalent to the
original formula at every successor. Formulas can be padded by
$A\dashv\vdash D_{u,y}G_yA$.

The formula selector $M_v(A,B)$ uses $C$ on corresponding sides of two
equations and distributes through aligned $D$ binders. If one formula
is an equation and the other has a $D$ binder, pad the equation using
the preceding gate before continuing. This finite recursive definition,
spelled out in Appendix~\ref{sec:choice-introduction}, gives
$M_0(A,B)\dashv\vdash A$ and $M_{Sv}(A,B)\dashv\vdash B$.
Write $Q_yA=D_{u,y}A$ with $u$ unused. One complete macro inventory is
\[
\begin{aligned}
 \bot_*&=(0=S0),\\
 A\to_*B&=D_{u,y}M_y(A,B),\\
 A\wedge_*B&=D_{u,y}G_y M_{p(y)}(A,B),\\
 \exists_*x\,A(x)&=D_{x,y}G_yA(x),\\
 A\vee_*B&=\exists_*x\,M_x(A,B),\\
 \forall_*x\,A(x)&=A(0)\wedge_* Q_y A(y).
\end{aligned}
\]
Every newly bound variable is fresh, except the explicitly displayed
argument of the formula being bound. These macros have all the usual
intuitionistic introduction and elimination rules. For example,
implication uses the selector's zero row as its assumption and its
successor rows as its conclusion; conjunction is read at the two controls
$1,2$; existential quantification uses the true zero gate. Universal
elimination uses $A(0)$ to activate the $Q$ clause. Under $0=S0$,
selector computation makes any two terms equal, and induction on finite
formula syntax yields explosion.

The appended inherited logical development proves these claims. Its
equivalent universal macro uses a gated paired conjunct; the macro displayed
here is the more direct $A(0)\wedge Q_yA(y)$ also used by the current
Lean development. Both have the ordinary universal reading.

\subsection{The numerical translations}

From scalar terms to iPRA, translate each fixed body $h(u,\vec p)$ by
structural recursion on its finite description. Its iterator is the
primitive-recursive operation
\[
 i_h(0,a,\vec p)=a,\qquad
 i_h(Sn,a,\vec p)=h(i_h(n,a,\vec p),\vec p).
\]
No number is asked to encode a function on which a universal total
evaluator operates.

In the other direction, the inherited scalar development first derives
predecessor from scalar iteration and equation induction. It then proves
successor injectivity, constructs an equality indicator that reflects
actual equality internally, and obtains induction on all quantifier-free
formulas. Pairing and its projections are derived next. For a fixed
primitive recursion
\[
 f(0,\vec p)=g(\vec p),\qquad
 f(Sn,\vec p)=h(n,f(n,\vec p),\vec p),
\]
the scalar state stores a pair consisting of the counter and the current
value. If $L,H$ are its two projections, set
\[
\begin{aligned}
 W(v,\vec p)&=\pi(SL(v),h^*(L(v),H(v),\vec p)),\\
 V(n,\vec p)&=I(n,\pi(0,g^*(\vec p));v.W(v,\vec p)),\\
 f^*(n,\vec p)&=H(V(n,\vec p)).
\end{aligned}
\]
Equation induction proves $L(V(n,\vec p))=n$, and the two recursion
equations for $f^*$ follow. The order of this argument matters: the
predecessor and equality-reflection proofs do not presuppose the general
primitive-recursion compiler. The appended foundations supply that order
and the necessary internal proofs~\cite{choice,bridge,erasure,forward,progress}.

\begin{theorem}[Required bi-interpretability and equiconsistency]
\label{thm:biinterp}
The ordinary realizations of $\mathcal E_3$ and $\mathcal E^0_3$ are
bi-interpretable with full first-order iPRA. Both numerical interpretation
domains are whole, equality is actual equality, and both numerical
composite isomorphisms have graph $x=y$. Both formula composites are
provably equivalent to the identity. Each presentation is equiconsistent
with iPRA.
\end{theorem}
\begin{proof}
Theorem~\ref{thm:raw-equivalence} transfers the full native consequence
relation to the raw uniform-clause reference, whose scalar and iPRA
interpretations are the ones just described. The zero companion has the
same exact comparison with its SeedEval reference.

For the numerical composites, induction on a fixed finite primitive-
recursive description shows that its double translation satisfies the
original defining equations. Internal equation induction proves uniqueness
of that fixed recursion. The analogous induction on scalar syntax proves
the other composite. These are statements for arbitrary internal counters,
not just for standard numeral instances.

The crucial native formula composite is
\[
 D_{x,y}B\dashv\vdash
 \exists_*x\bigl(B(x,0)\to_*\forall_*y\,B(x,y)\bigr).
\]
Forward, open its witness clause, assume its base, and use that clause
at a fresh variable before introducing the ordinary macros. Backward,
eliminate the existential at a fresh witness and use case introduction
with the supplied implication and universal at the successor instance.
The other formula composites follow from the selector and gate identities
and induction on finite formulas. The zero companion adds the internally
provable equivalence of a term with its zero-normal form.

The translations preserve contradiction: ordinary falsity maps to
$0=S0$, and the ordinary reading of $0=S0$ implies falsity by separation.
Effective translations of finite derivations in both directions therefore
give equiconsistency. This proof does not replace the whole domain by a
standard cut, a quotient, or a coded subset.
\end{proof}

The inherited numerical and logical bridge has a Lean formalization in
the pinned repository. The new fusion theorem and the association of its
Python checker with its stated rules are written proofs with executable
validation. No new Lean or Rocq verification is claimed.
